% ===== PRÉAMBULE CANONIQUE — à coller VERBATIM en tête de CHAQUE .tex =====
\documentclass[11pt,a4paper]{article}
\usepackage[utf8]{inputenc}\usepackage[T1]{fontenc}\usepackage{lmodern}
\usepackage[french]{babel}
\usepackage{amsmath,amssymb,mathtools}\usepackage{icomma}
\usepackage{siunitx}\sisetup{locale=FR}  % \qty{}{}, \si{}, \ang{}
\usepackage{xcolor}\usepackage{colortbl}
\usepackage{tikz}\usepackage{pgfplots}\pgfplotsset{compat=1.18}
\usepackage[siunitx,europeanresistors]{circuitikz} % circuits (élec.)
\usepackage[most]{tcolorbox}\usepackage{enumitem}\usepackage{array,booktabs}
\usepackage[a4paper,margin=2.2cm]{geometry}
\usetikzlibrary{arrows.meta,calc,angles,quotes,positioning,patterns,%
  backgrounds,shapes.geometric,decorations.pathmorphing,decorations.markings,intersections}
\definecolor{primary}{RGB}{222,49,99}\definecolor{primarydark}{RGB}{150,22,55}
\definecolor{defcol}{RGB}{0,114,178}\definecolor{methcol}{RGB}{52,92,148}
\definecolor{excol}{RGB}{0,135,90}\definecolor{attncol}{RGB}{200,40,55}
\tcbset{boxbase/.style={enhanced,breakable,boxrule=0.9pt,arc=2.5pt,
  left=3mm,right=3mm,top=2.3mm,bottom=2.3mm,fonttitle=\bfseries,coltitle=white}}
% --- boîtes TUTORIEL (noms EXACTS, ne pas renommer) ---
\newtcolorbox{cle}[1][]{boxbase,colback=defcol!6,colframe=defcol,
  title={\textsc{À retenir}\ifx\relax#1\relax\relax\else\textmd{ : \itshape #1}\fi}}
\newtcolorbox{methode}[1][]{boxbase,colback=methcol!7,colframe=methcol,sharp corners,
  title={\textsc{Méthode}\ifx\relax#1\relax\relax\else\textmd{ --- \itshape #1}\fi}}
\newtcolorbox{exemplebox}[1][]{boxbase,colback=excol!5,colframe=excol,
  title={\textsc{Exemple}\ifx\relax#1\relax\relax\else\textmd{ : \itshape #1}\fi}}
\newtcolorbox{piege}[1][]{boxbase,colback=attncol!6,colframe=attncol,
  title={\textsc{Piège}\ifx\relax#1\relax\relax\else\textmd{ : \itshape #1}\fi}}
\newtcolorbox{remarque}[1][]{boxbase,colback=black!4,colframe=black!45,
  title={\textsc{Remarque}\ifx\relax#1\relax\relax\else\textmd{ : \itshape #1}\fi}}
% --- boîtes EXERCICE / CORRIGÉ (noms EXACTS, auto-comptées) ---
\newtcolorbox[auto counter]{exo}[1][]{boxbase,colback=primary!4,colframe=primary,
  title={\textsc{Exercice}~\thetcbcounter\ifx\relax#1\relax\relax\else\textmd{ --- \itshape #1}\fi}}
\newtcolorbox[auto counter]{corrige}[1][]{boxbase,colback=excol!4,colframe=excol,
  title={\textsc{Corrigé}~\thetcbcounter\ifx\relax#1\relax\relax\else\textmd{ --- \itshape #1}\fi}}
\newcommand{\vv}[1]{\vec{#1}}\newcommand{\ds}{\displaystyle}
\newcommand{\R}{\mathbb{R}}\newcommand{\N}{\mathbb{N}}\newcommand{\Z}{\mathbb{Z}}
% (un \usepackage{fancyhdr}+en-tête est possible ; il est ignoré côté web)
% =========================================================================
\begin{document}
\sisetup{per-mode=symbol}

\begin{exo}[le travail du poids sur un tour complet]
Une personne de masse $m=\qty{70}{\kilogram}$ est assise dans une nacelle d'une grande roue de rayon
$R=\qty{20}{\metre}$. La roue effectue \emph{un tour complet} et la personne revient exactement à son
point de départ. On prend $g=\qty{10}{\metre\per\second\squared}$.
\begin{center}
\begin{tikzpicture}[>=Stealth,scale=0.8]
  \draw[black!60,line width=1pt] (0,0) circle (1.7);
  \foreach \a in {0,45,...,315} \draw[black!35] (0,0)--(\a:1.7);
  \fill[black!70] (0,0) circle (1.2pt);
  \draw[fill=primary!20,draw=black!70] (90:1.7)++(-0.18,-0.32) rectangle ++(0.36,0.32);
  \draw[->,line width=1pt,attncol] (90:1.7)++(0,-0.32)--++(0,-0.9) node[right,black]{\small $\vv G$};
  \draw[->,black!55] (35:2.0) arc (35:-35:2.0);
  \node at (2.55,0){\small tour};
  \draw[<->,black!45] (0,-2.0)--(0,0) node[midway,left]{\small $R$};
\end{tikzpicture}
\end{center}
\begin{enumerate}[label=\alph*)]
  \item En comparant l'altitude de départ et l'altitude d'arrivée, que vaut la variation d'énergie
        potentielle $\Delta E_p$ sur un tour complet ?
  \item En déduire le travail du poids $W(\vv G)$ sur ce tour, à l'aide de la relation
        $W(\vv G)=-\Delta E_p$. Le résultat dépend-il de $R$ ou de $m$ ?
\end{enumerate}
\end{exo}

\begin{exo}[monter un meuble : la rampe change la force, pas le travail]
Des déménageurs montent un meuble de masse $m=\qty{100}{\kilogram}$ jusqu'à une hauteur
$h=\qty{4}{\metre}$. Ils utilisent une rampe rectiligne \emph{sans frottement} de longueur
$L=\qty{5}{\metre}$. On prend $g=\qty{10}{\metre\per\second\squared}$.
\begin{center}
\begin{tikzpicture}[>=Stealth,scale=0.9]
  \draw[fill=primary!8,draw=black!70] (0,0)--(4,3)--(4,0)--cycle;
  \foreach \x in {0,0.4,...,3.8} \draw[black!45] (\x,0)--(\x-0.2,-0.22);
  \draw[fill=defcol!20,draw=black!70,rotate around={36.87:(2,1.5)}] (1.7,1.5) rectangle (2.5,2.0);
  \draw[<->,black!55,dashed] (4.5,0)--(4.5,3) node[midway,right]{\small $h=\qty{4}{\metre}$};
  \node at (1.4,1.55){\small $L=\qty{5}{\metre}$};
\end{tikzpicture}
\end{center}
\begin{enumerate}[label=\alph*)]
  \item Calcule le travail à fournir \emph{contre la pesanteur} pour élever le meuble de $h$.
        Ce travail dépend-il de la longueur de la rampe ?
  \item En l'absence de frottement, l'intensité $F$ de la force à exercer \emph{le long} de la rampe
        vérifie $F\times L=mgh$. Calcule $F$ et compare-la au poids $mg$. Quel est l'intérêt de la rampe ?
\end{enumerate}
\end{exo}

\begin{exo}[à quoi sert le travail du moteur d'une moto qui roule à vitesse constante ?]
Une moto roule \emph{à vitesse constante} sur une route rectiligne et \emph{horizontale}. Sur une
distance $d$, elle subit son poids $\vv G$, la réaction normale $\vv N$, la force motrice $\vv F_m$ et
la résultante des frottements $\vv F_f$.
\begin{enumerate}[label=\alph*)]
  \item La vitesse étant constante, que vaut la variation d'énergie cinétique $\Delta E_c$ entre le
        début et la fin du trajet ? Qu'impose alors le théorème de l'énergie cinétique à la somme des
        travaux $\sum W$ ?
  \item Sur cette route horizontale, que valent $W(\vv G)$ et $W(\vv N)$ ? Déduis-en la relation entre
        le travail de la force motrice $W(\vv F_m)$ et celui des frottements $W(\vv F_f)$.
\end{enumerate}
\end{exo}

\begin{exo}[puissance moyenne d'un moteur au démarrage]
Une voiture de masse $m=\qty{1000}{\kilogram}$ passe de l'arrêt à $v=\qty{72}{\kilo\metre\per\hour}$ en
$\Delta t=\qty{10}{\second}$. On néglige les frottements et on suppose le mouvement horizontal.
\begin{enumerate}[label=\alph*)]
  \item Convertis $v$ en \si{\metre\per\second}, puis calcule l'énergie cinétique acquise $E_c$.
  \item Comme $\Delta E_c$ est fourni par le travail du moteur, calcule le travail moteur $W$ puis la
        \textbf{puissance moyenne} $P=\dfrac{W}{\Delta t}$. Exprime-la en \si{\kilo\watt}.
\end{enumerate}
\end{exo}

\begin{exo}[vrai ou faux : le signe et la valeur d'un travail]
Indique si chaque affirmation est \textbf{vraie} ou \textbf{fausse}, en justifiant.
\begin{enumerate}[label=\alph*)]
  \item Le travail du poids d'un wagon qui roule de \qty{10}{\metre} sur des rails horizontaux est nul.
  \item En mouvement circulaire uniforme, le travail de la force centripète est toujours nul.
  \item Le travail d'une force de frottement de glissement est positif.
  \item Pour fournir une même quantité de travail, le moteur le plus puissant est celui qui met le
        \emph{plus} de temps.
\end{enumerate}
\end{exo}
\end{document}
